% META: {"id":"1SPE-DERLOCAL-CO-045","chapitre":"1SPE-DERIVATION-LOCAL","type_objet":"corrige","exercice_id":"1SPE-DERLOCAL-EX-045","status":"approved","origin":"generated","status_history":["generated","audited","accepted","approved"],"release_acceptance":"RELEASE_OWNER_BATCH_ACCEPTANCE","acceptance_closure_digest":"sha256:9b3ccf9a81c5520fbb7e03b7b2d3e7bf2057a02834b6d908bf3be5f49f3b553f","acceptance_receipt_digest":"sha256:4fad86f0282e0fa4e0419cca1ad6379ae7af5a280c04a4a90880f90a1c044242"}
% BEGIN-VERIFY
% from sympy import symbols, Rational
% x = symbols('x')
% f = (1 + x)**5
% assert f.subs(x, 0) == 1
% fp = f.diff(x)
% assert fp.subs(x, 0) == 5
% approx = 1 + 5*Rational(2, 100)
% assert approx == Rational(11, 10)
% END-VERIFY
\begin{corrige}{1SPE-DERLOCAL-EX-045}
\textbf{1.} On calcule :
\[
f(0) = (1 + 0)^5 = 1
\]
On admet $f'(x) = 5(1+x)^4$, donc :
\[
f'(0) = 5 \times (1+0)^4 = 5
\]

\textbf{2.} L'approximation affine de $f$ au voisinage de $x = 0$ s'écrit :
\[
f(x) \approx f(0) + f'(0) \cdot x = 1 + 5x
\]

\textbf{3.} On remarque que $(1{,}02)^5 = f(0{,}02)$. En appliquant l'approximation affine avec $x = 0{,}02$ :
\[
f(0{,}02) \approx 1 + 5 \times 0{,}02 = 1 + 0{,}1 = \boxed{1{,}1}
\]

\textbf{4.} La valeur exacte est $(1{,}02)^5 \approx 1{,}104\,1$. L'approximation affine donne $1{,}1$, soit une erreur d'environ $0{,}004$, c'est-à-dire moins de $0{,}4\,\%$ en erreur relative. L'approximation est très précise car $x = 0{,}02$ est proche de $0$.
\end{corrige}
